\chapter{SİSTEM TASARIMI VE GERÇEKLEŞTİRİM}
\label{IkinciBolum}

Bu bölümde, geliştirilen güvenli e-posta kimlik doğrulama sisteminin yazılım mimarisi, modüler yapısı, GUI gerçekleştirimi, animasyon modülleri ve test altyapısı ayrıntılı biçimde ele alınmaktadır.

\section{Sistem Mimarisi}

Uygulama, tek yönlü bağımlılık ilkesine göre tasarlanmış katmanlı bir mimariye sahiptir. Şekil \ref{fig:Mimari}'de görüldüğü üzere üç temel katman belirlenmiştir:

\begin{enumerate}
    \item \textbf{Kriptografi Katmanı (\texttt{crypto\_core.py}):} RSA anahtar üretimi, SHA-256 özetleme, RSA-PSS imzalama/doğrulama ve AES-256-GCM şifreleme/deşifreleme işlevlerini gerçekleştiren arka uç. GUI'den bağımsız; yalnızca Python standart kütüphanesi ve PyCA \texttt{cryptography} paketi kullanılmaktadır.
    \item \textbf{Sunum Katmanı (\texttt{alice\_panel.py}, \texttt{bob\_panel.py}):} Alice ve Bob panellerini yöneten, adım adım iş akışını tetikleyen ve animasyon pencerelerini çağıran bileşenler.
    \item \textbf{Animasyon Katmanı (\texttt{animation\_modals/}):} AES-256, RSA-2048 ve SHA-256 algoritmalarının iç işleyişini görselleştiren bağımsız pencere bileşenleri.
\end{enumerate}

\begin{table}[t]
	\caption{Proje Modülleri ve Sorumlulukları.}
	\centering
	\begin{tabular}{|l|l|}
		\hline \textbf{Modül} & \textbf{Sorumluluk} \\
		\hline \texttt{main\_gui.py} & Ana pencere, iş akışı yöneticisi \\
		\hline \texttt{crypto\_core.py} & Kriptografik iş mantığı (arka uç) \\
		\hline \texttt{alice\_panel.py} & Alice paneli ve animasyon konteyneri \\
		\hline \texttt{bob\_panel.py} & Bob paneli ve deşifre diyagramı \\
		\hline \texttt{theme.py} & Renk paleti ve tema sabitleri \\
		\hline \texttt{toast.py} & Doğrulama bildirimi bileşeni \\
		\hline \texttt{utils.py} & Yardımcı fonksiyonlar \\
		\hline \texttt{animation\_modals/aes\_animation.py} & AES-256-GCM adım animasyonu \\
		\hline \texttt{animation\_modals/rsa\_animation.py} & RSA-2048 eğitim animasyonu \\
		\hline \texttt{animation\_modals/sha256\_animation.py} & SHA-256 sıkıştırma animasyonu \\
		\hline \texttt{animation\_modals/aes\_pure.py} & Saf Python AES-256 (tur verisi) \\
		\hline \texttt{animation\_modals/sha256\_pure.py} & Saf Python SHA-256 (adım verisi) \\
		\hline
	\end{tabular}
	\label{tab:Moduller}
\end{table}

\section{Kriptografik İş Akışı}

Proje, \texttt{CryptoCore} sınıfı etrafında organize edilmiştir. RSA anahtar çifti \texttt{RSAKeyPair} veri sınıfında tutulmakta; şifreli paket ise \texttt{EncryptedPacket} veri sınıfıyla temsil edilmektedir. Temel veri sınıfları Tablo \ref{tab:VeriSiniflari}'da özetlenmiştir.

\begin{table}[t]
	\caption{Temel Veri Sınıfları.}
	\centering
	\begin{tabular}{|l|l|l|}
		\hline \textbf{Sınıf} & \textbf{Alan} & \textbf{Açıklama} \\
		\hline \texttt{RSAKeyPair} & \texttt{private\_key} & RSA-2048 gizli anahtar \\
		\hline                    & \texttt{public\_key}  & RSA-2048 açık anahtar \\
		\hline \texttt{EncryptedPacket} & \texttt{encrypted\_message} & AES-GCM şifreli metin + etiketi \\
		\hline                          & \texttt{encrypted\_session\_key} & OAEP şifreli $K_S$ \\
		\hline                          & \texttt{nonce} & 96-bit rastgele sayı \\
		\hline                          & \texttt{associated\_data} & AES-GCM AAD alanı \\
		\hline \texttt{StepResult} & \texttt{step\_number} & Adım numarası \\
		\hline                     & \texttt{step\_name}   & Adım adı \\
		\hline                     & \texttt{data}         & Ara çıktı sözlüğü \\
		\hline
	\end{tabular}
	\label{tab:VeriSiniflari}
\end{table}

Alice tarafının şifreleme sürecinde gerçekleştirilen altı adım Algoritma \ref{algo:Alice}'de özetlenmiştir.

\begin{algorithm}
	~\\
	\SetKwInOut{Input}{Girişler}
	\SetKwInOut{Output}{Çıkış}
	\Input{
		\textit{m}: Düz metin mesaj. \\
		$K^-_A$: Alice'in gizli RSA anahtarı. \\
		$K^+_B$: Bob'un açık RSA anahtarı. \\
	}
	\Output{\textit{Paket}: $\{C, K^+_B(K_S), \text{Nonce}\}$}
	\begin{algorithmic}[H]
		\STATE $H \leftarrow \text{SHA-256}(m)$
		\STATE $\sigma \leftarrow \text{RSA-PSS}_{K^-_A}(H)$
		\STATE $P \leftarrow m \| \sigma$
		\STATE $K_S \leftarrow \text{os.urandom}(32)$, $\text{Nonce} \leftarrow \text{os.urandom}(12)$
		\STATE $C \leftarrow \text{AES-256-GCM}_{K_S, \text{Nonce}}(P)$
		\STATE $E_{K_S} \leftarrow \text{RSA-OAEP}_{K^+_B}(K_S)$
		\STATE \textbf{return} $\{C, E_{K_S}, \text{Nonce}\}$
	\end{algorithmic}
	\caption{Alice Tarafı Şifreleme Algoritması}
	~\\
	\label{algo:Alice}
\end{algorithm}

\section{Animasyon Modülleri}

Proje, kriptografik algoritmaların iç işleyişini adım adım görselleştiren üç bağımsız animasyon penceresi içermektedir. Bu pencereler \texttt{CryptoAnimationWindow} temel sınıfından türetilmiş olup PyQt6'nın \texttt{QPainter} API'si kullanılarak özel grafik bileşenleri olarak gerçekleştirilmiştir.

\subsection{AES-256-GCM Animasyonu}

\texttt{aes\_animation.py} modülü, AES-256'nın 14 turunu görselleştirmektedir. Saf Python AES-256 gerçekleştirimi (\texttt{aes\_pure.py}) her tur sonrası 4×4 durum matrisini (\texttt{MatrixWidget}) depolar. Animasyon penceresi şu bileşenleri içermektedir:

\begin{itemize}
    \item Tıklanabilir tur çubuğu (Round 0 – Round 14).
    \item \textbf{SubBytes:} S-Box ile byte değiştirme, hücre hücre animasyon.
    \item \textbf{ShiftRows:} Satır kaydırma (0, 1, 2, 3 byte sola) için renkli ok gösterimi.
    \item \textbf{MixColumns:} GF($2^8$) sabit matris çarpımı formülü ($02 \cdot s_0 \oplus 03 \cdot s_1 \oplus s_2 \oplus s_3$).
    \item \textbf{AddRoundKey:} XOR ile tur anahtarı karıştırma.
    \item ECB ve GCM çıktılarının neden farklı olduğunu açıklayan eşleşme ekranı.
\end{itemize}

\subsection{RSA-2048 Animasyonu}

\texttt{rsa\_animation.py} modülü, RSA-2048'in matematiksel temellerini yedi adımda sunmaktadır: asal sayı seçimi, $n = p \cdot q$, $\phi(n)$, $e$ seçimi, $d$ hesabı (Genişletilmiş Öklid Algoritması adım adım), DER/Base64 dönüşümü ve gerçek 2048-bit anahtar gösterimi. Demo değerler ($p = 61$, $q = 53$) üzerinden şifreleme ve deşifreleme doğrulaması ($m = 65 \rightarrow c = 2790 \rightarrow m = 65$) gerçekleştirilmektedir.

\subsection{SHA-256 Animasyonu}

\texttt{sha256\_animation.py} modülü üç bölümden oluşmaktadır: dolgu görselleştirmesi, mesaj genişletme tablosu ($W_0 \ldots W_{31}$) ve QPainter tabanlı sıkıştırma devresi diyagramı. Sıkıştırma diyagramı; sekiz register (A–H), $T_1 = \Sigma_1(E) + \text{Ch}(E,F,G) + H + K_i + W_i$ ve $T_2 = \Sigma_0(A) + \text{Maj}(A,B,C)$ kutularını ve kaydırma bağlantılarını içermektedir. Her blok için 9 anlık görüntü (tur 1, 9, 17, 25, 33, 41, 49, 57, 64) sunulmaktadır.

\section{Test Altyapısı ve Sonuçları}

Uygulama, dört ayrı test dosyasında toplam 26 birim testi ile doğrulanmıştır. Testler Python \texttt{pytest} çerçevesi ile çalıştırılmaktadır.

\begin{table}[t]
	\caption{Test Dosyaları ve Kapsamları.}
	\centering
	\begin{tabular}{|l|c|l|}
		\hline \textbf{Test Dosyası} & \textbf{Test Sayısı} & \textbf{Kapsam} \\
		\hline \texttt{test\_crypto\_core.py} & 26 & SHA-256, RSA, AES-GCM, uçtan uca iş akışı \\
		\hline \texttt{test\_aes\_pure.py} & -- & Saf Python AES-256 gerçekleştirimi \\
		\hline \texttt{test\_sha256\_pure.py} & -- & Saf Python SHA-256 gerçekleştirimi \\
		\hline \texttt{test\_diagram\_rects.py} & -- & Diyagram dikdörtgen sınır kontrolleri \\
		\hline
	\end{tabular}
	\label{tab:Testler}
\end{table}

Temel test senaryoları arasında şunlar yer almaktadır:

\begin{itemize}
    \item SHA-256 çıktısının uzunluğunun ve bilinen test vektörlerinin doğrulanması.
    \item RSA anahtarlarının üretilmesi ve PEM/DER serileştirmesinin kontrolü.
    \item İmza doğrulamasının $H(m)$ üzerinde yapıldığının ve $H(H(m))$ ile yapılmadığının test edilmesi (\texttt{test\_signature\_is\_over\_hash\_not\_double\_hash}).
    \item AES-256-GCM şifreleme-deşifreleme döngüsünün bütünlüğünün doğrulanması.
    \item Bozuk şifreli metin veya bozuk imza ile gerçekleştirilen negatif test senaryoları.
    \item Alice'ten Bob'a uçtan uca iletim ve başarılı doğrulama testi.
\end{itemize}

\clearpage
